\subsection{Independent Multi-Objective Holdout}
The independent holdout compares the frozen DQL policy with the threshold-based
Heuristic and Static policies from Table~\ref{tab:policy-roles}. Full study
designs, training traces, distributions, and action-contract diagnostics are
retained in Section~\suppref{app:holdout-details}.

Table~\suppref{tab:locked-study-design} summarizes the evidence-generating runs.
In v3, three candidate networks were trained for exactly $10{,}000$
transitions under seeds 101, 211, and 307.  Candidate selection used ten fixed
validation seeds and required subscription-count rebalancing
($a^{\mathrm{reb}N}$), backlog rebalancing ($a^{\mathrm{reb}G}$), and full
connection migration ($a^{\mathrm{mig}}$) in at least $80\%$ of admissible
episodes in their intended scenarios.  Candidate 0 (seed 101) passed this gate
and was frozen before the holdout seed bank was read.

The independent v3 holdout contains ten previously unused seeds for each of
four scenarios and 500 decision intervals per episode.  The connection-count
skew scenario uses 1,000 subscriptions; the tiered-backlog scenario uses 1,200
subscriptions; and the client-slowdown and pod-pressure-skew scenarios each use
1,600 subscriptions.  Every scenario has four modeled pods, four initial
persistent connections, and a maximum of six connections.  The tiered-backlog
scenario assigns $25\%$ of its queues to each initial backlog level
$100{,}000$, $50{,}000$, $5{,}000$, and zero.  The client-slowdown scenario
reduces acknowledgement capacity to $35\%$ of its nominal value between
intervals 125 and 375.  The other two scenarios isolate initial
subscription-count skew and the placement of three of four connections on one
pod, respectively.

For policy $\pi$, paired episode $j$, and metric $v$, let
$Y_{\pi,j,v}$ denote the episode mean.  We define a direction-adjusted paired
difference
\begin{equation}
  D_{j,v}=d_v\left(Y_{\mathrm{DQL},j,v}
  -Y_{\mathrm{base},j,v}\right),
  \qquad
  d_v=\begin{cases}
  +1,&\text{higher is better},\\
  -1,&\text{lower is better}.
  \end{cases}
  \label{eq:direction-adjusted-effect}
\end{equation}
Thus, a positive $D_{j,v}$ always favors DQL.  The relative benefit reported
below is
\begin{equation}
  \mathrm{RB}_v=100\,
  \frac{\overline{D}_v}{\left|\overline{Y}_{\mathrm{base},v}\right|}.
  \label{eq:relative-benefit}
\end{equation}
Overbars in these two equations denote means across paired episodes.
The metric index $v$ is distinct from the client-population size $m$.
We resample the 40 paired differences $10{,}000$ times to obtain percentile
$95\%$ confidence intervals.  Two-sided paired sign-flip tests are exact for
the scenario-level sets of ten pairs and use $100{,}000$ Monte Carlo sign
vectors for the pooled set of 40 pairs.  Pooled $p$-values are adjusted by the
Holm procedure across the complete family of 18 metric-by-baseline tests.  The
scenario box plots in Figure~\suppref{fig:v3-scenario-reward} describe quartiles
and individual episode outcomes; the paired randomization tests
provide the statistical comparison.

We evaluate DQL as a sequential controller through return, operational
outcomes, paired effect sizes, and action coverage. Classifier accuracy and
ROC--AUC apply to labelled prediction tasks. Applying them here would mainly
measure agreement with a heuristic that supplied the correct-action labels,
rather than the controller's operating performance.

Across all 40 paired episodes, frozen DQL improves mean reward over the
heuristic from $-0.1609$ to $-0.1291$.  The mean paired difference is $0.0318$
(bootstrap 95\% CI $[0.0191,0.0455]$), equivalent to a $19.8\%$ relative
benefit under Eq.~\eqref{eq:relative-benefit}; the paired standardized effect is
$d_z=0.729$, and the Holm-adjusted sign-flip $p$-value is below $0.001$.
Table~\ref{tab:v3-paired-effects} and Figure~\suppref{fig:v3-paired-benefit} show
that this overall gain comes with tradeoffs.

\begin{table}[!htbp]
  \caption{Frozen DQL versus the heuristic on the independent v3 holdout.
  Positive benefit favors DQL.  CIs are paired-bootstrap intervals transformed
  by Eq.~\eqref{eq:relative-benefit}; $p_H$ is Holm-adjusted over 18 pooled
  metric-by-baseline tests.}
  \Description{Table listing DQL and heuristic means, direction-adjusted
  relative benefits with 95 percent confidence intervals, and Holm-adjusted
  paired sign-flip p-values for nine operational metrics.}
  \label{tab:v3-paired-effects}
  \centering
  \small
  \begin{tabular}{lrrrr}
    \toprule
    Metric & DQL & Heuristic & Benefit [95\% CI] (\%) & $p_H$ \\
    \midrule
    Mean reward & $-0.1291$ & $-0.1609$ & $19.8\ [11.9,28.3]$ & $<.001$ \\
    Backlog ratio & $0.1749$ & $0.2002$ & $12.7\ [-1.2,27.1]$ & $.095$ \\
    Conn.-subscription variance & $574.7$ & $5768.7$ & $90.0\ [61.0,120.8]$ & $<.001$ \\
    Conn.-backlog variance & $2.280\!\times\!10^{13}$ & $1.568\!\times\!10^{13}$ & $-45.4\ [-72.5,-22.6]$ & $.022$ \\
    Subscription imbalance & $0.0898$ & $0.4343$ & $79.3\ [59.1,99.5]$ & $<.001$ \\
    Backlog imbalance & $0.2194$ & $0.1638$ & $-34.0\ [-56.4,-13.6]$ & $.022$ \\
    Combined conn. imbalance & $0.1546$ & $0.2991$ & $48.3\ [33.8,63.8]$ & $<.001$ \\
    Pod pressure & $0.1878$ & $0.2051$ & $8.4\ [3.6,13.6]$ & $.022$ \\
    Throughput ratio & $0.2840$ & $0.2973$ & $-4.5\ [-7.1,-2.2]$ & $.022$ \\
    \bottomrule
  \end{tabular}
\end{table}

The largest benefit is a $90.0\%$ reduction in connection-level subscription
variance, consistent with the design objective of moving subscriptions away
from overloaded persistent connections.  Subscription imbalance falls by
$79.3\%$, combined connection imbalance by $48.3\%$, and mean pod pressure by
$8.4\%$.  Mean backlog ratio is $12.7\%$ lower, but its interval crosses zero
and the adjusted test is not significant ($p_H=.095$).  Conversely, DQL has
$45.4\%$ higher connection-backlog variance, $34.0\%$ higher normalized backlog
imbalance, and $4.5\%$ lower throughput than the heuristic.  The scalar reward
still favors DQL because its configured penalties place substantial weight on
subscription pressure and combined imbalance. The holdout therefore identifies
a multi-objective advantage for DQL and a queue-drainage advantage for the
heuristic.

\subsection{Adverse-Probe Ablation}
The separate 30-seed, 1,200-subscription ablation compares the same frozen DQL
policy with and without one forced promotion to the highest-pressure probe
Pod. Guarded DQL reduces normalized backlog AUC from 0.68938 to 0.66046:
a paired difference of $-0.028918$ (95\% CI $[-0.039444,-0.019073]$), or 4.19\%.
The guarded heuristic retains the lower AUC of 0.59104. This isolates the
benefit of guarding promotion, not universal superiority of learned selection.
The complete five-policy comparison, methods, and separate five-seed
20,000-subscription confirmation are retained in
Section~\suppref{sec:adverse-probe-ablation}.
\paragraph{Reviewer artifact access.}
The accompanying supplementary archive, \texttt{reviewer-evidence.zip},
contains episode-level simulator CSVs, declared protocols, original aggregate
analyses, the frozen 21--32--6 checkpoint, and identifier-free live summaries.
Its README gives the file map and exact offline commands. A Python 3.9+
standard-library script verifies the checksum manifest, checks paired cohort
coverage, recomputes policy means and selected effect sizes, and reconciles
the loaded trial's reported subscription and ACK totals. A separate inference
option evaluates the saved network without transport access or model updates.
The archive supports reanalysis and checkpoint inspection; it does not expose
corporate credentials, raw event payloads, private endpoint identifiers, or
the proprietary service source. It is not a complete independent deployment
rerun. Provenance records distinguish unchanged measurements from redacted
local-path metadata.

The public research repository is available at
\url{https://github.com/manish94u/rebalancePersistance}.
It provides the manuscript sources, standalone simulator, redacted evidence,
frozen checkpoint, and step-by-step replication instructions.
